\documentclass{report}
\usepackage{amsmath,amssymb}
\setlength{\parindent}{0mm}
\setlength{\parskip}{1em}
\begin{document}
\begin{center}
\rule{6in}{1pt} \
{\large Andrew Reisner \\
{\bf Multilevel Solvers for High Resolution Electric Field Calculations}}

107 N Busey Apt B \\ Urbana \\ IL 61801
\\
{\tt areisne2@illinois.edu}\\
Luke Olson\end{center}

High fidelity electric field calculations are a critical component in plasma
simulations. In this this talk we consider the problem of a dielectric
barrier
discharge (DBD) wherein the electric field is calculated to support a
compressible flow, thus requiring a highly efficient global solve. The
electric
field is constructed on a logically structured but mapped mesh which yields
anisotropy in the operator, along with jumps in the permittivity. Another
challenge arises in the modeling of a dielectric barrier discharge, where
dielectrics result in localized Dirichlet blocks within the domain. In a
multilevel solver, these interior blocks are not resolved on coarse grids,
leading to a deterioration in convergence with a strong dependence on the
alignment and size relative to the coarse levels. We investigate the
dependence of various multilevel solvers in this context and in a parallel
setting. In particular, we detail the convergence of multilevel methods for
high resolution electric field calculations in the presence of warped meshes
with jumping coefficients.


\end{document}
